%---------------------------------------------------------------------
% UNIT 9 --- OPTIMIZATION
% Source: Section 9.9 Self Assessment Questions, pp. 205-206.
%---------------------------------------------------------------------
\chapter{Optimization}

\srcnote{Source: \emph{Business Mathematics} 5405/1429, \S9.9 \saq{Self Assessment Questions}, pp.~205--206.}

\begin{keyidea}[What the two derivatives tell you]
\begin{center}
\small\sffamily
\begin{tabular}{@{}L{3.4cm} L{4.6cm} L{5.4cm}@{}}
\toprule
 & \textbf{$f'$ --- first derivative} & \textbf{$f''$ --- second derivative} \\
\midrule
Zero at        & stationary points        & possible inflection points \\
Positive means & $f$ increasing           & $f$ concave up (holds water) \\
Negative means & $f$ decreasing           & $f$ concave down (spills water) \\
\bottomrule
\end{tabular}
\end{center}

\medskip
\textbf{First derivative test:} $f'$ changing $+\to-$ gives a maximum,
$-\to+$ a minimum, no change gives neither.

\textbf{Second derivative test:} at a stationary point,
$f''<0\Rightarrow$ maximum, $f''>0\Rightarrow$ minimum,
$f''=0\Rightarrow$ inconclusive.

\medskip
\textbf{An inflection point needs $f''$ to change sign}, not merely to vanish.
\end{keyidea}

\begin{keyidea}[The four business functions]
\[
R(x)=x\cdot p(x),\qquad
P(x)=R(x)-C(x),\qquad
\bar{C}(x)=\frac{C(x)}{x},
\]
and their marginals are the derivatives $MR=R'$, $MC=C'$, $MP=P'$. Profit is
maximised where $P'=0$, i.e.\ where $MR=MC$, provided $P''<0$.
\end{keyidea}

%---------------------------------------------------------------------
\section{Increase, decrease, concavity and inflection}

\begin{question}{Q.(a)}
Find the intervals on which $f$ is increasing, decreasing, concave up and
concave down, and locate any points of inflection.
\textbf{(i)} $f(x)=x^{2}-3x+8$ \quad
\textbf{(ii)} $f(x)=\dfrac{x^{4}}{3}-\dfrac{x}{3}$ \quad
\textbf{(iii)} $f(x)=\dfrac{x}{x^{2}+2}$ \quad
\textbf{(iv)} $f(x)=x^{4}-5x^{3}+9x^{2}$
\end{question}

\begin{solution}
\textbf{(i)} $f'(x)=2x-3$, zero at $x=\tfrac32$. Since $f'<0$ to the left and
$f'>0$ to the right:
\[
\text{decreasing on }\left(-\infty,\tfrac32\right),\qquad
\text{increasing on }\left(\tfrac32,\infty\right).
\]
$f''(x)=2>0$ everywhere, so the curve is concave up throughout and has
\emph{no} inflection point. The stationary point $x=\tfrac32$ is a minimum,
$f(\tfrac32)=\tfrac94-\tfrac92+8=\tfrac{23}{4}=5.75$.

\medskip
\textbf{(ii)} $f(x)=\tfrac13(x^{4}-x)$, so
\[
f'(x)=\frac{4x^{3}-1}{3}=0 \;\Longrightarrow\; x^{3}=\frac14
\;\Longrightarrow\; x=4^{-1/3}\approx 0.630 .
\]
\[
\text{decreasing on }(-\infty,\,0.630),\qquad
\text{increasing on }(0.630,\,\infty).
\]
\[
f''(x)=\frac{12x^{2}}{3}=4x^{2}\ge 0 \quad\text{for all }x .
\]
$f''$ vanishes at $x=0$ but does \emph{not} change sign there, so the curve is
concave up everywhere and has no inflection point.

\medskip
\textbf{(iii)} Quotient rule:
\[
f'(x)=\frac{(x^{2}+2)-x(2x)}{(x^{2}+2)^{2}}=\frac{2-x^{2}}{(x^{2}+2)^{2}} .
\]
The denominator is always positive, so the sign of $f'$ is the sign of
$2-x^{2}$: zero at $x=\pm\sqrt2$.
\[
\text{increasing on }(-\sqrt2,\sqrt2),\qquad
\text{decreasing on }(-\infty,-\sqrt2)\text{ and }(\sqrt2,\infty).
\]
Differentiating again and simplifying:
\[
f''(x)=\frac{2x\left(x^{2}-6\right)}{(x^{2}+2)^{3}} .
\]
Zeros at $x=0$ and $x=\pm\sqrt6$; the denominator is positive, so the sign is
that of $2x(x^{2}-6)$:

\begin{center}
\small\sffamily
\begin{tabular}{@{}C{3.2cm} C{2.0cm} C{2.4cm} C{3.0cm}@{}}
\toprule
Interval & $2x$ & $x^{2}-6$ & $f''$ \\
\midrule
$x<-\sqrt6$          & $-$ & $+$ & $-$ \ concave down \\
$-\sqrt6<x<0$        & $-$ & $-$ & $+$ \ concave up \\
$0<x<\sqrt6$         & $+$ & $-$ & $-$ \ concave down \\
$x>\sqrt6$           & $+$ & $+$ & $+$ \ concave up \\
\bottomrule
\end{tabular}
\end{center}

$f''$ changes sign at all three zeros, so there are inflection points at
$x=-\sqrt6$, $x=0$ and $x=\sqrt6$, with
$f(0)=0$ and $f(\pm\sqrt6)=\pm\tfrac{\sqrt6}{8}$.

\medskip
\textbf{(iv)} $f'(x)=4x^{3}-15x^{2}+18x=x\left(4x^{2}-15x+18\right)$.

The quadratic factor has discriminant $225-288=-63<0$ and positive leading
coefficient, so it is positive for every $x$ and never contributes a root. Hence
$f'=0$ only at $x=0$, and the sign of $f'$ is the sign of $x$:
\[
\text{decreasing on }(-\infty,0),\qquad \text{increasing on }(0,\infty),
\]
with a relative minimum at $x=0$, $f(0)=0$.
\[
f''(x)=12x^{2}-30x+18=6\left(2x^{2}-5x+3\right)=6(2x-3)(x-1),
\]
zero at $x=1$ and $x=\tfrac32$:
\[
\text{concave up on }(-\infty,1)\text{ and }\left(\tfrac32,\infty\right);
\qquad
\text{concave down on }\left(1,\tfrac32\right).
\]
Both are genuine inflection points, at $(1,\,5)$ and
$\left(\tfrac32,\,\tfrac{135}{16}\right)=(1.5,\,8.4375)$.

\ans{(i) decreasing then increasing about $x=\tfrac32$; concave up everywhere;
no inflection. \quad
(ii) decreasing then increasing about $x=4^{-1/3}\approx0.630$; concave up
everywhere; no inflection. \quad
(iii) increasing on $(-\sqrt2,\sqrt2)$, decreasing outside; inflections at
$x=-\sqrt6,\,0,\,\sqrt6$. \quad
(iv) decreasing on $(-\infty,0)$, increasing on $(0,\infty)$; concave down only
on $\left(1,\tfrac32\right)$; inflections at $(1,5)$ and $(1.5,\,8.4375)$.}
\end{solution}

\begin{pitfall}
In (ii), $f''(0)=0$ but $f''=4x^{2}$ is non-negative on both sides, so there is
\emph{no} inflection at $x=0$. ``$f''=0$'' identifies a candidate; only a sign
change confirms it.
\end{pitfall}

%---------------------------------------------------------------------
\section{Relative extrema by both derivative tests}

\begin{question}{Q.(b)}
Use both the first and the second derivative test to find the relative extrema
of $f$.
\textbf{(i)} $f(x)=x^{4}-2x^{3}+4x^{2}$ \quad
\textbf{(ii)} $f(x)=x^{2}(x-2)^{2}$ \quad
\textbf{(iii)} $f(x)=3x+2x^{1/3}$ \quad
\textbf{(iv)} $f(x)=\dfrac{x^{2}}{x^{4}+x^{2}-1}$ \quad
\textbf{(v)} $f(x)=\dfrac{x+4}{x-3}$
\end{question}

\begin{solution}
\textbf{(i)} $f'(x)=4x^{3}-6x^{2}+8x=2x\left(2x^{2}-3x+4\right)$.

The quadratic has discriminant $9-32=-23<0$, so it is always positive and
contributes no roots. Thus $x=0$ is the only critical point.

\emph{First derivative test:} the sign of $f'$ is the sign of $2x$, so $f'$
changes $-\to+$ at $x=0$: a relative minimum.

\emph{Second derivative test:} $f''(x)=12x^{2}-12x+8$, and $f''(0)=8>0$:
minimum. Both agree. $f(0)=0$.

\medskip
\textbf{(ii)} Recognise $f=\left[x(x-2)\right]^{2}=(x^{2}-2x)^{2}$, then use the
chain rule:
\[
f'(x)=2(x^{2}-2x)(2x-2)=4x(x-2)(x-1).
\]
Critical points $x=0,\,1,\,2$.

\emph{First derivative test:}
\begin{center}
\small\sffamily
\begin{tabular}{@{}C{2.6cm} C{1.6cm} C{1.8cm} C{1.8cm} C{2.0cm} C{2.6cm}@{}}
\toprule
Interval & $4x$ & $(x-2)$ & $(x-1)$ & $f'$ & $f$ \\
\midrule
$x<0$     & $-$ & $-$ & $-$ & $-$ & decreasing \\
$0<x<1$   & $+$ & $-$ & $-$ & $+$ & increasing \\
$1<x<2$   & $+$ & $-$ & $+$ & $-$ & decreasing \\
$x>2$     & $+$ & $+$ & $+$ & $+$ & increasing \\
\bottomrule
\end{tabular}
\end{center}
So $x=0$ is a minimum, $x=1$ a maximum and $x=2$ a minimum.

\emph{Second derivative test:} $f''(x)=4\left(3x^{2}-6x+2\right)$, so
\[
f''(0)=8>0 \ (\text{min}),\qquad
f''(1)=-4<0 \ (\text{max}),\qquad
f''(2)=8>0 \ (\text{min}).
\]
Values: $f(0)=0$, $f(1)=1$, $f(2)=0$. The two tests agree throughout.

\medskip
\textbf{(iii)} $f(x)=3x+2x^{1/3}$, so
\[
f'(x)=3+\frac{2}{3}x^{-2/3}=3+\frac{2}{3x^{2/3}} .
\]
Since $x^{2/3}>0$ for every $x\neq0$, $f'>3>0$ everywhere it is defined, and
$f'$ is undefined (infinite) at $x=0$. There is no point where $f'=0$, and $f'$
does not change sign across $x=0$.

Therefore $f$ is strictly increasing on its whole domain and has \textbf{no
relative extrema}. The second derivative test is not applicable: it requires a
stationary point, and there is none.

\medskip
\textbf{(iv)} $f(x)=\dfrac{x^{2}}{x^{4}+x^{2}-1}$. By the quotient rule,
\[
f'(x)=\frac{2x(x^{4}+x^{2}-1)-x^{2}(4x^{3}+2x)}{(x^{4}+x^{2}-1)^{2}}
=\frac{2x^{5}+2x^{3}-2x-4x^{5}-2x^{3}}{(x^{4}+x^{2}-1)^{2}}
\]
\[
=\frac{-2x^{5}-2x}{(x^{4}+x^{2}-1)^{2}}
=\frac{-2x\left(x^{4}+1\right)}{(x^{4}+x^{2}-1)^{2}} .
\]
The denominator is positive and $x^{4}+1>0$ always, so the sign of $f'$ is the
sign of $-2x$: positive for $x<0$, negative for $x>0$.

\emph{First derivative test:} $f'$ changes $+\to-$ at $x=0$, so there is a
relative \textbf{maximum} at $x=0$, with
\[
f(0)=\frac{0}{0+0-1}=0 .
\]
The domain excludes the points where $x^{4}+x^{2}-1=0$: putting $u=x^{2}$,
$u^{2}+u-1=0$ gives $u=\tfrac{-1+\sqrt5}{2}\approx0.618$, so
$x\approx\pm0.786$. Between those two vertical asymptotes the denominator is
negative, so $f\le0$ there --- consistent with $f(0)=0$ being the largest value
on that interval.

\medskip
\textbf{(v)} $f(x)=\dfrac{x+4}{x-3}$:
\[
f'(x)=\frac{(x-3)-(x+4)}{(x-3)^{2}}=\frac{-7}{(x-3)^{2}} .
\]
This is negative for every $x\neq3$ and never zero. There are no critical
points, so $f$ has \textbf{no relative extrema}; it is strictly decreasing on
$(-\infty,3)$ and on $(3,\infty)$ separately.
\[
f''(x)=\frac{14}{(x-3)^{3}}:\quad
\text{concave down for }x<3,\quad\text{concave up for }x>3 .
\]
There is a vertical asymptote at $x=3$ and a horizontal asymptote at $y=1$.

\ans{(i) Relative minimum at $(0,0)$. \quad
(ii) Minima at $(0,0)$ and $(2,0)$, maximum at $(1,1)$. \quad
(iii) None --- $f$ is strictly increasing. \quad
(iv) Relative maximum at $(0,0)$. \quad
(v) None --- $f$ is strictly decreasing on each branch.}
\end{solution}

\begin{pitfall}
In (iii) and (v), ``no relative extrema'' is the full answer. A function with
no stationary point has nothing to classify, and the second derivative test
cannot be applied. Do not invent an extremum at a point where the derivative is
merely undefined ((iii), $x=0$) or where the function itself is undefined
((v), $x=3$).
\end{pitfall}

%---------------------------------------------------------------------
\section{Curve sketching}

\begin{question}{Q.(c)}
Sketch the graph of the functions
\textbf{(i)} $f(x)=2x^{3}-3x^{2}-12x+1$ \quad
\textbf{(ii)} $f(x)=\dfrac{3x^{2}}{x^{2}+5}$ \quad
\textbf{(iii)} $f(x)=x-\dfrac1x$ \quad
\textbf{(iv)} $f(x)=\dfrac{x-1}{x+1}$ \quad
\textbf{(v)} $f(x)=x^{3}$
\end{question}

\begin{solution}
A sketch is a report on four things: stationary points, concavity,
asymptotes, and intercepts. Take (i) and (ii) in full.

\medskip
\textbf{(i) $f(x)=2x^{3}-3x^{2}-12x+1$.}
\[
f'(x)=6x^{2}-6x-12=6(x^{2}-x-2)=6(x-2)(x+1),
\]
so stationary points at $x=-1$ and $x=2$:
\[
f(-1)=-2-3+12+1=8,\qquad f(2)=16-12-24+1=-19 .
\]
\[
f''(x)=12x-6,\qquad f''(-1)=-18<0 \ (\text{max}),\qquad f''(2)=18>0 \ (\text{min}).
\]
$f''=0$ at $x=\tfrac12$, where it changes sign: an inflection point at
$\left(\tfrac12,\,-5.5\right)$. So: increasing, then decreasing between $-1$ and
$2$, then increasing again; concave down for $x<\tfrac12$, concave up after.

\medskip
\textbf{(ii) $f(x)=\dfrac{3x^{2}}{x^{2}+5}$.}
\[
f'(x)=\frac{6x(x^{2}+5)-3x^{2}(2x)}{(x^{2}+5)^{2}}=\frac{30x}{(x^{2}+5)^{2}} ,
\]
zero only at $x=0$, changing $-\to+$: a minimum at $(0,0)$.
\[
f''(x)=\frac{150-90x^{2}}{(x^{2}+5)^{3}},
\]
zero at $x=\pm\sqrt{5/3}\approx\pm1.291$, giving inflection points at
$\left(\pm1.291,\;0.75\right)$. As $x\to\pm\infty$, $f\to3$: a horizontal
asymptote at $y=3$, approached from below. There is no vertical asymptote,
since $x^{2}+5$ never vanishes.

\medskip
\textbf{(iii)} $f(x)=x-\tfrac1x$ has $f'=1+\tfrac{1}{x^{2}}>0$ always, so it is
increasing on each branch, with a vertical asymptote at $x=0$ and the slant
asymptote $y=x$.

\textbf{(iv)} $f(x)=\tfrac{x-1}{x+1}$ has $f'=\tfrac{2}{(x+1)^{2}}>0$, so it is
increasing on each branch, with a vertical asymptote at $x=-1$, a horizontal
asymptote at $y=1$, and intercepts $(1,0)$ and $(0,-1)$.

\textbf{(v)} $f(x)=x^{3}$ has $f'=3x^{2}\ge0$ (increasing everywhere, with a
horizontal tangent but no extremum at $x=0$) and $f''=6x$, giving an inflection
point at the origin.

\begin{figure}[H]
\centering
\begin{tikzpicture}
\begin{axis}[
  width=11.5cm, height=7cm,
  axis lines=middle, xlabel={$x$}, ylabel={$f(x)$},
  xmin=-3.2, xmax=4.6, ymin=-26, ymax=18,
  xtick={-3,-2,...,4}, ytick={-25,-20,...,15},
  grid=both, grid style={line width=0.2pt, draw=bmgrey!25},
  tick label style={font=\footnotesize}, label style={font=\small},
  every axis plot/.append style={very thick}
]
\addplot[bmblue, domain=-2.6:4.1, samples=120] {2*x^3 - 3*x^2 - 12*x + 1};
\addplot[only marks, mark=*, mark size=2.4pt, bmred]
  coordinates {(-1,8) (2,-19) (0.5,-5.5)};
\node[font=\scriptsize\sffamily, anchor=south west, bmred] at (axis cs:-1.0,8.4)
  {max $(-1,\,8)$};
\node[font=\scriptsize\sffamily, anchor=north west, bmred] at (axis cs:2.1,-19.4)
  {min $(2,\,-19)$};
\node[font=\scriptsize\sffamily, anchor=south west, bmred] at (axis cs:0.7,-5.2)
  {inflection $(0.5,\,-5.5)$};
\end{axis}
\end{tikzpicture}
\caption{Q.(c)(i): $f(x)=2x^{3}-3x^{2}-12x+1$, showing the maximum, the
minimum, and the inflection point where concavity reverses.}
\end{figure}

\begin{figure}[H]
\centering
\begin{tikzpicture}
\begin{axis}[
  width=11.5cm, height=5.6cm,
  axis lines=middle, xlabel={$x$}, ylabel={$f(x)$},
  xmin=-8, xmax=8, ymin=-0.5, ymax=3.6,
  grid=both, grid style={line width=0.2pt, draw=bmgrey!25},
  tick label style={font=\footnotesize}, label style={font=\small},
  every axis plot/.append style={very thick}
]
\addplot[bmgreen, domain=-8:8, samples=160] {3*x^2/(x^2+5)};
\addplot[bmgrey, dashed, domain=-8:8, thin] {3};
\addplot[only marks, mark=*, mark size=2.2pt, bmred]
  coordinates {(0,0) (1.2910,0.75) (-1.2910,0.75)};
\node[font=\scriptsize\sffamily, anchor=south east, bmgrey] at (axis cs:7.9,3.02)
  {$y=3$ asymptote};
\node[font=\scriptsize\sffamily, anchor=north west, bmred] at (axis cs:0.2,0.35)
  {min $(0,0)$};
\end{axis}
\end{tikzpicture}
\caption{Q.(c)(ii): $f(x)=3x^{2}/(x^{2}+5)$, with its minimum at the origin,
inflection points at $x=\pm\sqrt{5/3}$, and the horizontal asymptote $y=3$.}
\end{figure}

\ans{(i) Maximum $(-1,8)$, minimum $(2,-19)$, inflection $(0.5,-5.5)$. \quad
(ii) Minimum $(0,0)$, inflections at $(\pm1.291,\,0.75)$, horizontal asymptote
$y=3$. \quad
(iii) Increasing on each branch; asymptotes $x=0$ and $y=x$. \quad
(iv) Increasing on each branch; asymptotes $x=-1$ and $y=1$. \quad
(v) Increasing everywhere; inflection at the origin.}
\end{solution}

%---------------------------------------------------------------------
\section{Where profit is increasing}

\begin{question}{Q.(d)}
A manufacturer sells a product with cost and revenue functions
\[
C(x)=2.2x-0.0002x^{2},\qquad R(x)=4.5x-0.004x^{2},\qquad 0\le x\le 1200,
\]
where $x$ is the number of items sold. Find the interval on which the profit
function is increasing.
\end{question}

\begin{solution}
Build the profit function first, and collect like terms carefully --- the
subtraction of a negative $x^{2}$ term is where errors enter:
\[
\begin{aligned}
P(x)&=R(x)-C(x)\\
&=\left(4.5x-0.004x^{2}\right)-\left(2.2x-0.0002x^{2}\right)\\
&=4.5x-0.004x^{2}-2.2x+0.0002x^{2}\\
&=2.3x-0.0038x^{2} .
\end{aligned}
\]

Profit is increasing where $P'(x)>0$:
\[
P'(x)=2.3-0.0076x>0
\;\Longleftrightarrow\;
x<\frac{2.3}{0.0076}=302.63 .
\]

Since the domain is $0\le x\le 1200$:
\[
\text{$P$ increasing on } [0,\;302.63), \qquad
\text{$P$ decreasing on } (302.63,\;1200] .
\]
$P''=-0.0076<0$, so the turning point is a maximum:
\[
P(302.63)=2.3(302.63)-0.0038(302.63)^{2}=696.05-348.03=348.03 .
\]

\medskip
In whole units, output of about 303 items maximises profit at roughly \$348.

\ans{Profit is increasing on $0\le x<302.63$, i.e.\ up to about
\textbf{303 units}, after which it declines. Maximum profit $\approx\$348.03$.}
\end{solution}

\begin{pitfall}
$-(-0.0002x^{2})=+0.0002x^{2}$. Dropping that sign gives
$P=2.3x-0.0042x^{2}$ and a turning point at $274$ instead of $303$. Write out
the subtraction line in full rather than doing it in your head.
\end{pitfall}

%---------------------------------------------------------------------
\section{When is marginal revenue increasing?}

\begin{question}{Q.(e)}
The demand function is $p(x)=300-2x$. Determine when marginal revenue is
increasing.
\end{question}

\begin{solution}
Revenue is price times quantity:
\[
R(x)=x\,p(x)=x(300-2x)=300x-2x^{2}.
\]
Marginal revenue is the first derivative:
\[
MR(x)=R'(x)=300-4x .
\]
Marginal revenue is increasing where \emph{its} derivative is positive, i.e.\
where $R''>0$:
\[
\dd{}{x}\bigl(MR\bigr)=R''(x)=-4 .
\]
This is a negative constant --- it is never positive. Therefore marginal revenue
is \textbf{never increasing}: it is a decreasing linear function of $x$,
falling by \$4 for each extra unit sold, on the whole domain.

\medskip
\textbf{What is true instead.} $MR>0$ (revenue itself still rising) for
$x<75$; $MR=0$ at $x=75$, where revenue peaks at
$R(75)=300(75)-2(75)^{2}=22500-11250=11{,}250$; and $MR<0$ beyond that, where
selling more actually reduces revenue.

\ans{Marginal revenue $MR=300-4x$ has $R''=-4<0$ everywhere, so marginal revenue
is \textbf{never increasing} --- it decreases at a constant rate of 4 per unit.
Revenue itself increases only while $MR>0$, that is for $x<75$.}
\end{solution}

\begin{pitfall}
``Revenue increasing'' and ``marginal revenue increasing'' are different
questions. The first asks about $R'>0$; the second about $R''>0$. This question
asks the second, and the answer --- ``never'' --- is a legitimate one.
\end{pitfall}

%---------------------------------------------------------------------
\section{Minimising average cost}

\begin{question}{Q.(f)}
The total cost $c$ of producing $x$ units of a product is
$c=0.03x^{2}+8x+300$. At what level of output will the average cost per unit be
minimum?
\end{question}

\begin{solution}
Average cost is total cost divided by output:
\[
\bar{c}(x)=\frac{0.03x^{2}+8x+300}{x}=0.03x+8+\frac{300}{x}
=0.03x+8+300x^{-1}.
\]
Differentiate and set to zero:
\[
\bar{c}\,'(x)=0.03-\frac{300}{x^{2}}=0
\;\Longrightarrow\;
x^{2}=\frac{300}{0.03}=10{,}000
\;\Longrightarrow\;
x=100
\]
(taking the positive root, since output cannot be negative).

Second derivative test:
\[
\bar{c}\,''(x)=\frac{600}{x^{3}},\qquad \bar{c}\,''(100)=\frac{600}{10^{6}}>0
\;\Longrightarrow\;\text{minimum}.
\]
\[
\bar{c}(100)=0.03(100)+8+\frac{300}{100}=3+8+3=14 .
\]

\vcheck{$\bar{c}(90)=2.7+8+3.333=14.03$ and $\bar{c}(110)=3.3+8+2.727=14.03$.
Both exceed 14, confirming the minimum.}

\medskip
\textbf{Worth noting.} At the minimum, marginal cost equals average cost:
$c'(x)=0.06x+8$, so $c'(100)=14=\bar{c}(100)$. That identity always holds at
the minimum of average cost, and makes a good check.

\ans{Average cost is minimised at an output of \textbf{100 units}, where the
average cost is \textbf{\$14} per unit.}
\end{solution}

%---------------------------------------------------------------------
\section{Maximising profit from demand and cost}

\begin{question}{Q.(g)}
The demand $p$ and cost function $c$ for a product are
\[
p=52-0.02x,\qquad c(x)=400+40x .
\]
At what level of output will profit be maximised? At what price does this
profit arise, and what is the profit?
\end{question}

\begin{solution}
Revenue is price times quantity:
\[
R(x)=x\,p=x(52-0.02x)=52x-0.02x^{2}.
\]
Profit:
\[
P(x)=R(x)-c(x)=52x-0.02x^{2}-400-40x=12x-0.02x^{2}-400 .
\]
\[
P'(x)=12-0.04x=0 \;\Longrightarrow\; x=\frac{12}{0.04}=300 .
\]
\[
P''(x)=-0.04<0 \;\Longrightarrow\; \text{maximum}.
\]

Price at that output:
\[
p=52-0.02(300)=52-6=\$46 .
\]
Profit:
\[
P(300)=12(300)-0.02(300)^{2}-400=3600-1800-400=\$1400 .
\]

\vcheck{Via $MR=MC$: $MR=52-0.04x$ and $MC=40$, so $52-0.04x=40$ gives
$x=300$ --- the same output, reached by the economic rule rather than the
calculus one. And $R(300)=15600-1800=13{,}800$, $c(300)=400+12000=12{,}400$,
difference $=1{,}400$.}

\ans{Profit is maximised at an output of \textbf{300 units}, at a price of
\textbf{\$46} per unit, giving a maximum profit of \textbf{\$1{,}400}.}
\end{solution}

\begin{pitfall}
Maximum profit is not maximum revenue. Revenue here peaks at $x=1300$
(where $52-0.04x=0$), far beyond the profit-maximising 300. Optimise the
function the question names.
\end{pitfall}

%---------------------------------------------------------------------
\section{Pricing a cable subscription}

\begin{question}{Q.(h)}
A TV cable company has 500 subscribers paying \$20 per month. There will be 200
more subscribers for each \$0.60 decrease in the monthly fee. At what rate will
maximum revenue be obtained, and what will this revenue be?
\end{question}

\begin{solution}
Let $n$ be the number of \$0.60 reductions applied. Then

\begin{center}
\small\sffamily
\begin{tabular}{@{}L{4.6cm} L{7.4cm}@{}}
\toprule
Monthly fee & $20-0.6n$ dollars \\
Number of subscribers & $500+200n$ \\
\bottomrule
\end{tabular}
\end{center}

Revenue is fee times subscribers:
\[
\begin{aligned}
R(n)&=(20-0.6n)(500+200n)\\
&=10000+4000n-300n-120n^{2}\\
&=10000+3700n-120n^{2} .
\end{aligned}
\]
\[
R'(n)=3700-240n=0 \;\Longrightarrow\; n=\frac{3700}{240}=15.4167 .
\]
\[
R''(n)=-240<0 \;\Longrightarrow\;\text{maximum}.
\]

At that $n$:
\[
\text{fee}=20-0.6(15.4167)=20-9.25=\$10.75,
\]
\[
\text{subscribers}=500+200(15.4167)=3583.3,
\]
\[
R=10000+3700(15.4167)-120(15.4167)^{2}
=10000+57{,}041.67-28{,}520.83=\$38{,}520.83 .
\]

\medskip
\textbf{The whole-number answer.} $n$ must be a whole number of \$0.60
reductions, and subscribers come in whole people. Testing the two neighbours:

\begin{center}
\small\sffamily
\begin{tabular}{@{}C{1.6cm} C{2.6cm} C{2.8cm} C{3.0cm}@{}}
\toprule
$n$ & Fee & Subscribers & Revenue \\
\midrule
15 & \$11.00 & 3{,}500 & \$38{,}500 \\
16 & \$10.40 & 3{,}700 & \$38{,}480 \\
\bottomrule
\end{tabular}
\end{center}

So the best attainable plan is $n=15$: a fee of \$11.00 per month.

\ans{The continuous optimum is a monthly rate of \textbf{\$10.75} with
$\approx3{,}583$ subscribers, giving revenue of \textbf{\$38{,}520.83}.
Restricted to whole \$0.60 reductions, the best rate is \textbf{\$11.00} per
month with 3{,}500 subscribers, giving \textbf{\$38{,}500} per month.}
\end{solution}

\begin{pitfall}
Define your variable explicitly --- ``$n$ = number of \$0.60 decreases'' --- and
write both the fee and the subscriber count in terms of it before multiplying.
Setting up $R$ in terms of the fee directly is possible but invites sign errors,
and most lost marks in this question are lost before any calculus begins.
\end{pitfall}
